\documentclass{article}
\usepackage[T1]{fontenc}
\usepackage{lmodern}
\usepackage{graphicx}
\usepackage{amsmath,amssymb}
\usepackage[a4paper,margin=2cm]{geometry}
\usepackage[absolute,overlay]{textpos}
\setlength{\TPHorizModule}{1mm}
\setlength{\TPVertModule}{1mm}
\usepackage[indonesian]{babel}
\usepackage{hyperref}
\setlength{\emergencystretch}{2em}
% unit: Halaman 88

\begin{document}

% === Halaman 88 (llm) ===

\textbf{3. Cocok untuk perangkat dengan sumber daya terbatas}\\ Karena operasi internalnya sederhana, Salsa20 dapat diimplementasikan dengan baik pada berbagai perangkat, termasuk perangkat yang kemampuan komputasinya terbatas.\\
\vspace{5mm}
\textbf{4. Cocok untuk implementasi perangkat lunak}\\ Salsa20 dapat memperoleh performa yang baik hanya dengan operasi integer dasar. Karena itu Salsa20 sangat populer dalam komunitas kriptografi, dan desainnya kemudian menjadi dasar bagi \textbf{ChaCha}, yang akhirnya lebih banyak digunakan dalam protokol modern.\\
\vspace{5mm}
\textbf{5. Menghasilkan blok \textit{keystream} yang besar}\\ Satu blok Salsa20 menghasilkan $16 \times 32 \text{ bit} = 512 \text{ bit} = 64 \text{ byte keystream}$. Jika plaintext panjang, counter tinggal dinaikkan untuk menghasilkan blok berikutnya.

\end{document}
